\section{Introduction}
\label{sec:intro}

Continuous water-quality monitoring underpins the management of rivers that supply drinking water and receive pollution. In-situ sensors now measure temperature, specific conductance, dissolved oxygen, and pH every few minutes, and the resulting high-frequency records reveal diel cycles, storm responses, and short pulses that grab samples miss \citep{rode2016}. Wireless sensor networks have made such monitoring affordable \citep{adumanu2017}, and deployments in Ghana and Kosovo have shown that it works in practice, but they relied on cellular backhaul \citep{adumanu2020,ahmedi2022}. The rivers under the greatest pressure are often the furthest from that infrastructure; in Ghana, for example, small-scale gold mining has degraded rivers in forested catchments where terrestrial coverage is sparse \citep{duncan2020}. This motivating case is not evaluated here.

Satellite Internet of Things (IoT) services now reach such sites with small, low-power terminals \citep{centenaro2021}, but their capacity is sold in messages, not bandwidth. A typical module sends one-way messages with a payload of 23 bytes and receives no acknowledgement \citep{kineis2023kim1}. A four-parameter sensor sampling every 10 minutes produces 576 readings a day; even packed at 13 bits each, without timestamps, into the 20 bytes that remain after a three-byte message header, they would fill at least 47 messages. A monitoring agency therefore needs to know what a credible record of a river costs, how to deliver it when messages are scarce and may be lost without notice, and what the delivered record can be trusted for. Many uses of such records tolerate a delay of a day: compliance assessment, trend analysis, and model calibration work from daily or longer summaries, whereas alarms that require action within minutes are outside the scope of this study.

This study defined credibility by established practice: a record is sufficiently accurate when every reconstructed reading lies within the tolerance that the USGS requires of an excellent continuous record \citep{wagner2006}. Because the link returns nothing, the server must be able to establish from the bytes it received which readings meet that tolerance. The study asked four research questions:
\begin{enumerate}
  \item RQ1: How many values and messages does a record within tolerance require per day, and how does that need vary across rivers, seasons, and sampling intervals?
  \item RQ2: Which delivery design maximises the share of days that the server can fully certify over a one-way link, without knowledge of the loss process, compared with the designs used in practice?
  \item RQ3: Which message plans suffice, and can an adequate plan be chosen from a short pilot record?
  \item RQ4: Do the delivered records support the decisions that agencies take, such as detecting low dissolved oxygen or conductance pulses?
\end{enumerate}

The algorithms used are established: minimum-size piecewise-linear approximation within a maximum error \citep{imai1988}, budgeted compression of sensor history \citep{deligiannakis2007}, and erasure codes \citep{reed1960}. Satellite-linked animal tags have long sent compressed profiles pseudo-randomly over a one-way link \citep{boehme2009,photopoulou2015}. What had not been established is how these elements must be adapted and combined so that a server without feedback can certify an environmental record, what such a record costs on real rivers, and which decisions it supports. The contributions of this study are the following:
\begin{itemize}
  \item A framework that delivers error-bounded river records over one-way, message-priced links: a quantisation-exact, self-contained encoder; a byte-level message frame; a record-first erasure code; and a server-side certificate that is provably sound and needs no feedback.
  \item A proof that the record-first code maximises the probability of full certification among unit-level linear designs for every loss process, such that it needs no channel knowledge, with exact delivery probabilities that agree with byte-level simulation.
  \item An evaluation on \nSitesDev{} development sites and \nSitesConf{} sealed confirmation sites in five other regions, with the software and analysis plan frozen before the confirmation data were opened.
  \item A decision-oriented assessment: certified records never gave a wrong certain decision and detected short events that periodic and tag telemetry missed, but they do not resolve features finer than the tolerance.
\end{itemize}

Section~\ref{sec:related} reviews related work. Section~\ref{sec:model} describes the framework. Section~\ref{sec:study} describes the data and the study design. Section~\ref{sec:results} reports the results, Section~\ref{sec:discussion} discusses them, and Section~\ref{sec:conclusion} concludes.
